\scanpage{10}
% archive sheet number: 9
% Vertical emphasis stroke at left of the corollary.

forment un syst. projectif \inserted{strict} d'algèbres sur $A$. $\unsure{C_i|\mathcal C_i}(\unsure{G|B})_f$ est équivalente \unsure{canoniqt|canoniqmt|canoniquement} à la catégorie des $U_i$-modules \inserted{de t.f.}. Et les \deleted{foncteurs} \unsure{transitions|transition} \deleted{\unsure{entre}} les $C_i(G)$ sont isomorphes aux \deleted{\unsure{aux}} \deleted{foncteurs} \uline{restriction des scalaires}

\margin{\rule{.04em}{2.6em}}\uline{Corollaire} Sous les conditions précédentes, tt module \inserted{de t.f. $E$} d'une représentation de $G$ est quotient d'un module \inserted{\uline{projectif}} de t.f. où $G$ opère.

\uline{Prop} Soit $G$ gpe affine plat sur $A$ anneau de Dedekind, d'algèbre $B$ satisfaisant aux conditions du lemme 1 (\unsure{NB.} si $G$ est de t.f. cela signifie que pour tte composante irréd. de $G$, \uline{ttes} ses fibres sont non vides).

Soit \deleted{\unsure{$A'$ une $A$-algèbre}} $\unsure{S|S'}$ un schéma sur $S$, tel que sur $S$, tt module qu. coh. de t.f. \deleted{admette une \unsure{résolut.}} \inserted{soit isomorphe} à un quotient d'un module loc. libre de t.f.

\deleted{Alors} (p. ex. $S$ admet \unsure{modules|module} inv. \unsure{amples|ample}, ou $S$ noeth. régulier). Alors tt
